%% Self-contained captions. No cross-references, no section numbers: each
%% caption must be readable on its own if the figure is lifted out.

\newcommand{\CapCapacity}{%
Two ways of spending the same parameters. A subject-driven personalisation run
teaches a concept to a latent diffusion backbone through low-rank adapters on
the attention projections. The upper bar trains one adapter across both
attention types at once. The lower bar trains two adapters separately, one on
the cross-attention projections and one on the self-attention projections, and
switches both on at generation time. The two adapters together train
\CrossParams{} and \SelfParams{} parameters, which sum to the \JointParams{}
that the single adapter trains; the split is a partition of the same parameter
set rather than a change of capacity. Both arms were trained at
\TrainSteps{} steps, learning rate \LearningRate{}, adapter rank \LoraRank{},
resolution \Resolution{} and seed \TrainSeed{}. Counts are read from the
training records written by the runs.%
}

\newcommand{\CapCollapse}{%
The values a dispersion endpoint throws away. Each line follows one background
prompt across three arms of the same personalisation study: a single adapter
trained across both attention types, two separately trained adapters composed
at generation time, and the same composition with the self-attention adapter
switched off. All three arms were generated over the same \NPairs{} prompts, so
each line is a like-for-like sequence. The number printed beneath each arm is
the endpoint the recorded comparison used: the standard deviation of that arm's
\NPairs{} identity fidelities, computed across prompts. Collapsing four values
into one leaves a single number per arm, and a comparison of two such numbers
has nothing left to pair.%
}

\newcommand{\CapDispersion}{%
The recorded comparison, carried through to the interval it supports. Each row
is a ratio of dispersions: the standard deviation of identity fidelity across
\NPairs{} background prompts for one arm, divided by the same quantity for the
arm trained as a single adapter. Points are the observed ratios and bars are
95\% confidence intervals from an $F$ distribution at \NPairs{} generations per
arm, which is the interval an unpaired comparison of two dispersions at this
size admits. The dashed line marks equal dispersion. Both intervals contain it.
The interval assumes the two arms are independent; they are not, because they
share the background prompts, and that assumption is the one the unpaired
comparison already makes. The horizontal axis is logarithmic.%
}

\newcommand{\CapPaired}{%
The same generations, analysed as the pairs they were run as. Each point is one
background prompt, and its position is the composed arm's value minus the value
of the arm trained as a single adapter at that same prompt, for the endpoint
named above the panel. Identity fidelity is the quantity the recorded
comparison aggregated into a dispersion; prompt fidelity is recorded in the
same files and was not used. The grey band is the 95\% confidence interval for
the mean paired difference and the vertical rule inside it is that mean. The
annotation gives the interval and the exact two-sided sign-test probability on
the same pairs. With \NPairs{} pairs that probability cannot fall below
\FloorAtFour{}, so neither panel can reach a conventional threshold by the sign
test whatever the points do.%
}

\newcommand{\CapFloor}{%
What a sign test can reach, before any data. Each curve gives the smallest
two-sided probability an exact sign test can return at a given number of pairs:
the lower curve when every pair falls the same way, the others when one or two
pairs fall the other way. The quantity depends on the number of pairs alone, so
it is fixed before an experiment is run and no effect size changes it. The
dashed line is \Alpha{}. The comparison re-analysed here has \NPairs{} pairs
and so cannot cross that line; \MinNAny{} pairs is the first size that can, and
then only on unanimity. The pre-registered design uses \DesignPerConcept{}
pairs per concept, which reaches \Alpha{} while allowing \DesignTolerance{}
pairs to disagree.%
}

\newcommand{\CapArms}{%
The three adapters on disk, as their training records describe them. Trainable
parameter counts, wall clock and final training loss are read from the manifest
each run wrote. The two adapters in the lower block are the ones composed at
generation time; their parameter counts sum to the count of the adapter trained
across both attention types, which is what makes the two arms matched on
capacity. All three were trained at the same budget, learning rate, rank,
resolution and seed.%
}

\newcommand{\CapPairs}{%
Per-prompt values behind the recorded comparison, and the differences the
recorded endpoint discarded. Rows are the background prompts, joined across
arms on the prompt text. \emph{Joint} is the arm trained as one adapter across
both attention types; \emph{composed} is the arm whose two adapters were
trained separately and switched on together. Differences are composed minus
joint, so a negative value means the composed arm scored lower at that prompt.
The final row is the endpoint the recorded comparison used: the standard
deviation of each arm's four identity fidelities across prompts, which is one
number per arm and therefore admits no paired difference.%
}

\newcommand{\CapProtocol}{%
The paired experiment, fixed before it is run. The pairing unit, the arms, the
endpoints and the tests are stated here so that no choice among them can be
made after the results are seen. The number of background prompts is the
recorded grid, unchanged. The number of generation seeds is the smallest that
leaves each concept able to reach \Alpha{} on an exact sign test with at least
one pair falling the wrong way; every other size in the table follows from it.
The training budget is that of the arms whose comparison this protocol
re-tests. Nothing in this table has been executed.%
}
